\documentclass[11pt,a4paper]{article}

\usepackage[utf8]{inputenc}
\usepackage[T1]{fontenc}
\usepackage{amsmath,amssymb,amsthm}
\usepackage{booktabs}
\usepackage{graphicx}
\usepackage[margin=1in]{geometry}
\usepackage{hyperref}
\usepackage{natbib}
% microtype omitted: requires scalable fonts absent from a minimal TeX install

\hypersetup{colorlinks=true,linkcolor=blue,citecolor=blue,urlcolor=blue}

\newtheorem{theorem}{Theorem}
\newtheorem{proposition}[theorem]{Proposition}
\newtheorem{lemma}[theorem]{Lemma}
\theoremstyle{definition}
\newtheorem{definition}[theorem]{Definition}
\newtheorem{assumption}[theorem]{Assumption}
\theoremstyle{remark}
\newtheorem{remark}[theorem]{Remark}

\newcommand{\indep}{\perp\!\!\!\perp}
\newcommand{\Ecal}{\mathcal{E}}
\newcommand{\Scal}{\mathcal{S}}
\newcommand{\Gstar}{\mathcal{G}^{\star}}
\newcommand{\pa}{\mathrm{pa}}

\title{\bf Certificate-Only Causal Discovery:\\
Anytime-Valid Adjacency and Orientation}

\author{
  Quang-Vinh Dang\thanks{British University Vietnam, Hung Yen, Vietnam.
  Email: \texttt{vinh.dq4@buv.edu.vn}. ORCID: 0000-0002-3877-8024.}
}

\date{}

\begin{document}
\maketitle

\begin{abstract}
Constraint-based causal discovery removes an edge when a conditional-independence
test \emph{fails to reject}, converting absence of evidence into evidence of
absence. This is why finite-sample correctness for such methods requires
faithfulness, and it is why they cannot say which of their outputs the data
actually supports. We propose a discovery algorithm built on a single rule:
assert a feature of the graph only by \emph{rejecting} a null that is true
whenever that feature is absent, and never assert anything by failing to
reject. Two certificates implement the rule. For adjacency, the null ``some
conditioning set separates the pair'' is a union null, and the minimum of the
per-set e-processes is an e-process for it; crossing certifies an edge, using
no faithfulness assumption. For orientation, the two directions of a pair
induce different joint densities under a linear non-Gaussian model, so a
sequential universal-inference e-process against the reversed factorisation
certifies an arrowhead. Both are anytime-valid, so the estimate may be
monitored continuously and the run stopped on any data-dependent rule. The
resulting algorithm grows from the empty graph, reports undecided pairs as
undecided, and bounds the probability that any certified edge or arrowhead is
ever wrong, uniformly over stopping times. Empirically, on instances violating
$\delta$-strong faithfulness the algorithm's error stays within budget while
delete-on-non-rejection methods fail in 60--80\% of runs; and 96\% of true
edges are oriented against a Markov-equivalence-class ceiling of 33\%, because
orientation error control and orientation beyond the equivalence class have
previously been available only separately.
\end{abstract}

\noindent\textbf{Keywords:} causal discovery; e-values; anytime-valid inference;
safe testing; family-wise error rate; LiNGAM

\section{Introduction}
\label{sec:intro}

A causal discovery algorithm returns a graph, and a practitioner reading that
graph would like to know which of its features the data support. Constraint-based
methods---PC, FCI and their descendants \citep{spirtes2000causation}---answer
this question poorly, for a structural reason. They begin from the complete
graph and delete an edge whenever a conditional-independence test fails to
reject. Failure to reject is not evidence of independence: it is compatible
with independence, with a weak dependence, and with a test that lacked power.
The deletion nevertheless happens, and the resulting absence of an edge is
reported with the same authority as its presence.

Two consequences follow. First, finite-sample guarantees for these methods
require faithfulness, or a quantitative strengthening of it; the assumption is
doing the work that the data cannot. Second, the output cannot distinguish
``this edge is absent'' from ``nothing was learned about this pair''.

This paper takes the opposite discipline. We adopt a single rule:

\begin{quote}
\emph{Assert a feature of the graph only by rejecting a null hypothesis that is
true whenever that feature is absent. Never assert anything by failing to
reject.}
\end{quote}

The rule is restrictive, and its consequences are not cosmetic. An algorithm
obeying it cannot begin from the complete graph, because it has no licence to
delete; it begins from the empty graph and adds. It cannot report an edge as
absent, only as undecided. And what it controls is the probability of asserting
something false, rather than the probability of missing something true.

Making the rule operational requires two certificates, and the second has no
direct precedent.

\paragraph{Adjacency.} For a pair $(i,j)$ the natural null is composite:
\emph{some} conditioning set separates them. This is true whenever the pair is
non-adjacent, so rejecting it certifies an edge. Because the null is a union,
the minimum of the per-set e-processes is an e-process for it
(Section~\ref{sec:adjacency}). Two things follow that matter more than the
construction. Validity uses no faithfulness assumption---if $i$ and $j$ are
non-adjacent then $\pa(i)$ or $\pa(j)$ separates them, so the null holds---and
multiplicity runs over pairs rather than over (pair, conditioning set) triples.

\paragraph{Orientation.} Orientation error is not uncontrolled in the
literature: PC-p \citep{strobl2017estimating} tests unshielded colliders and
Meek-rule orientations and controls their false discovery rate. But those tests
are confined to the Markov equivalence class, and an edge lying in no
v-structure is unorientable in principle by any of them. Conversely the
functional-model methods that escape the equivalence class---LiNGAM and its
descendants \citep{shimizu2006linear,shimizu2011directlingam}---choose a
direction by comparing scores, with no error control on the choice. Under a
linear non-Gaussian model the two orientations of a pair induce different joint
densities and exactly one is correct, so ``the direction is $j \to i$'' is a
genuine null; a sequential universal-inference e-process against it certifies
the arrowhead $i \to j$ (Section~\ref{sec:direction}). Under Gaussian noise the
two factorisations are observationally indistinguishable, the process does not
grow, and no certificate is issued---the procedure abstains exactly where the
direction is unidentified.

\paragraph{Contributions.}
\begin{enumerate}
\item An e-process for the composite ``separability'' null obtained as a
  minimum over conditioning sets, giving time-uniform certification of
  adjacency under the Markov condition without faithfulness
  (Section~\ref{sec:adjacency}).
\item A sequential certificate for causal \emph{orientation} that is valid at
  any stopping time and is not confined to the Markov equivalence class,
  abstaining automatically under Gaussianity (Section~\ref{sec:direction}).
\item CERT-CD, an algorithm combining both under one budget, whose output is a
  three-valued graph in which every asserted edge and arrowhead carries a
  certificate and undecided pairs are explicit, with a whole-graph guarantee
  that is proved rather than assumed (Section~\ref{sec:algorithm}).
\item An empirical study separating instances by whether the premise of the
  competing guarantee holds, which shows that near-unfaithful instances cost
  the method power rather than validity (Section~\ref{sec:experiments}).
\end{enumerate}

We are explicit throughout about what is not new. Section~\ref{sec:related}
gives that accounting in detail.

\section{Background}
\label{sec:background}

\subsection{E-processes and Ville's inequality}

An \emph{e-variable} for a null $H_0$ is a non-negative statistic with
$\mathbb{E}_P[E] \le 1$ for all $P \in H_0$. An \emph{e-process} is a
non-negative process $(E_t)_{t \ge 0}$ such that $\mathbb{E}_P[E_\tau] \le 1$
for every stopping time $\tau$ and every $P \in H_0$
\citep{ramdas2023gametheoretic,grunwald2024safe}. Ville's inequality then gives
\begin{equation}
\label{eq:ville}
P\Bigl( \exists\, t \ge 0 : E_t \ge 1/\alpha \Bigr) \;\le\; \alpha ,
\end{equation}
a bound holding \emph{uniformly in time}. This is what licenses continuous
monitoring: the analyst may inspect the evidence after every observation and
stop whenever they wish, without inflating the error rate. A fixed-sample
p-value offers no such licence, and the gap is large in practice
(Section~\ref{sec:exp-calibration}).

\subsection{Causal discovery setting}

Let $\Gstar$ be a directed acyclic graph over variables $X_1,\dots,X_d$ and let
$P$ be Markov to $\Gstar$. Write $\Scal_k(i,j)$ for the set of all conditioning
sets $S \subseteq V \setminus \{i,j\}$ with $|S| \le k$; this family is fixed
before any data are seen. We say a pair is \emph{separable within $\Scal_k$} if
$X_i \indep X_j \mid X_S$ for some $S \in \Scal_k$.

\begin{assumption}[Separator size]
\label{ass:sep}
Every non-adjacent pair of $\Gstar$ has a separating set of size at most $k$.
\end{assumption}

\noindent Assumption~\ref{ass:sep} is structural rather than distributional and
is satisfied whenever $k$ is at least the maximum in-degree, since $\pa(i)$
d-separates $i$ from any non-adjacent $j$ that is not a descendant of $i$, and
$\pa(j)$ does so otherwise.

\section{The adjacency certificate}
\label{sec:adjacency}

Fix a pair $(i,j)$ and consider the composite null
\begin{equation}
\label{eq:sepnull}
H^{\mathrm{sep}}_{ij} \;:\quad
\exists\, S \in \Scal_k(i,j) \ \text{ such that } \ X_i \indep X_j \mid X_S .
\end{equation}
By the Markov property and Assumption~\ref{ass:sep}, $H^{\mathrm{sep}}_{ij}$
holds whenever $i$ and $j$ are non-adjacent in $\Gstar$. Rejecting it therefore
certifies an edge. Let $(E^{(S)}_t)$ be any e-process for the point null
$X_i \indep X_j \mid X_S$, and define
\begin{equation}
\label{eq:adjacency}
A_t(i,j) \;=\; \min_{S \in \Scal_k(i,j)} E^{(S)}_t(i,j) .
\end{equation}

\begin{proposition}
\label{prop:adjacency}
$(A_t)$ is an e-process for $H^{\mathrm{sep}}_{ij}$.
\end{proposition}

\begin{proof}
Let $P \in H^{\mathrm{sep}}_{ij}$, so $X_i \indep X_j \mid X_{S^\star}$ under
$P$ for some $S^\star \in \Scal_k(i,j)$. Then $A_t \le E^{(S^\star)}_t$
pointwise, and $E^{(S^\star)}$ is an e-process for a null that $P$ satisfies, so
for every stopping time $\tau$,
$\mathbb{E}_P[A_\tau] \le \mathbb{E}_P[E^{(S^\star)}_\tau] \le 1$.
\end{proof}

That a minimum of e-values is an e-value for a union null is standard
\citep{vovk2021evalues}; what matters here is the identification of adjacency
as the complement of such a union, and the consequences below.

\begin{remark}[No faithfulness]
\label{rem:faithfulness}
The proof uses only that the null~\eqref{eq:sepnull} is true for non-adjacent
pairs, which follows from the Markov condition and
Assumption~\ref{ass:sep}. Faithfulness is needed only for \emph{power}:
specifically adjacency-faithfulness \citep{ramsey2006adjacency}, which is what
makes every $E^{(S)}$ grow for a genuinely adjacent pair. An instance where it
fails leaves pairs undecided rather than producing wrong edges. Stated exactly,
validity requires three things---the Markov condition, Assumption~\ref{ass:sep},
and a valid per-set e-process. The third is not free: our default primitive is
a linear-Gaussian safe test (Section~\ref{sec:primitive}), so that model is
assumed there, and substituting a nonparametric sequential test
\citep{he2026sequential} removes it. Proposition~\ref{prop:adjacency} is
indifferent to which primitive is used.
\end{remark}

\begin{remark}[Multiplicity over pairs]
One hypothesis is tested per pair, so a union bound runs over $\binom{d}{2}$
hypotheses rather than $\binom{d}{2} \cdot |\Scal_k|$. At $d = 6$, $k = 2$ this
is a $24\times$ larger per-hypothesis level than a triple-indexed budget.
\end{remark}

\subsection{The per-set primitive}
\label{sec:primitive}

For linear-Gaussian data we use an exact e-process obtained from a Bayes factor
with right-Haar nuisance priors. Writing the hypothesis in regression form
$X_i = X_S^\top \alpha + \beta X_j + \varepsilon$ with
$\varepsilon \sim N(0,\sigma^2)$, the null is $\beta = 0$ with nuisance
$(\alpha,\sigma)$. Placing $\pi(\alpha,\sigma) \propto \sigma^{-1}$ on the
nuisance block under both hypotheses and $\beta \mid \sigma \sim N(0,\sigma^2 g)$
on the tested coefficient, and integrating in closed form, gives
\begin{equation}
\label{eq:safelinear}
E_n \;=\; (1+G)^{-1/2}\,\bigl(1 - \kappa\,\hat r_n^2\bigr)^{-(n-p)/2},
\qquad
G = g\,S_{yy}, \quad \kappa = \tfrac{G}{1+G},
\end{equation}
where $p = |S| + 1$, $S_{xx},S_{xy},S_{yy}$ are cross-products of the
$X_S$-residualised vectors and $\hat r_n^2 = S_{xy}^2/(S_{xx}S_{yy})$ is the
squared sample partial correlation. Two properties drive the implementation.
First, $E_n$ depends on the data only through Schur complements of the running
Gram matrix, so a single $O(d^2)$ sufficient statistic serves the entire
hypothesis family and no per-hypothesis state is stored. Second, $g$ must be
fixed in advance---a prior scale adapting to the running sample size destroys
the martingale property---so we set it from a warm-up prefix excluded from every
e-process.

Calibration was verified by simulation rather than asserted: monitoring after
every observation over 3000 replications and 1500 time points gave realised
rates of $0.149$, $0.076$, $0.041$ and $0.008$ against nominal
$\alpha \in \{0.20, 0.10, 0.05, 0.01\}$.

Existing sequential conditional-independence tests largely operate under the
Model-X framework, which requires the conditional law of $X_i$ given $X_S$;
this is typically inaccessible in discovery, an obstacle also noted by
\citet{csillag2025prediction}. Construction~\eqref{eq:safelinear} avoids it.

\section{The direction certificate}
\label{sec:direction}

Under a linear non-Gaussian acyclic model \citep{shimizu2006linear} the two
orientations of a pair induce different joint densities and exactly one is
correct. Hence
\begin{equation}
H^{\mathrm{dir}}_{j \to i} \;:\quad
\text{the pair was generated by the } j \to i \text{ factorisation}
\end{equation}
is a null that is true whenever the arrowhead $i \to j$ is absent, and rejecting
it certifies that arrowhead. Conditioning on a covariate block $Z$, the two
competing factorisations are
\begin{align}
D^{j \to i} &: \quad X_j = a^\top Z + u, \qquad X_i = \beta X_j + \gamma^\top Z + e, \\
D^{i \to j} &: \quad X_i = b^\top Z + v, \qquad X_j = \delta X_i + \eta^\top Z + w,
\end{align}
both being conditional densities of $(X_i, X_j)$ given $Z$, so the marginal law
of $Z$ cancels. We test by sequential universal inference
\citep{wasserman2020universal}:
\begin{equation}
\label{eq:direction}
\log E_t \;=\; \sum_{s \le t} \log D^{i \to j}_{\hat\theta_{s-1}}(o_s)
\;-\; \sup_{\theta} \sum_{s \le t} \log D^{j \to i}_{\theta}(o_s),
\end{equation}
with the numerator's parameters fitted on strictly past data.

\begin{proposition}
\label{prop:direction}
$(E_t)$ defined by~\eqref{eq:direction} is an e-process for
$H^{\mathrm{dir}}_{j \to i}$.
\end{proposition}

\begin{proof}
Fix the true null parameter $\theta^\star$ and set
$M_t = \prod_{s \le t} D^{i \to j}_{\hat\theta_{s-1}}(o_s) /
D^{j \to i}_{\theta^\star}(o_s)$. Since $\hat\theta_{s-1}$ is
$\mathcal{F}_{s-1}$-measurable and $D^{i \to j}_{\hat\theta_{s-1}}$ is a
probability density, $\mathbb{E}[M_s \mid \mathcal{F}_{s-1}] = M_{s-1}$, so $M$
is a non-negative martingale with $M_0 = 1$. The denominator
of~\eqref{eq:direction} is a supremum over the null model and therefore at
least the likelihood at $\theta^\star$, giving $E_t \le M_t$ pointwise.
Optional stopping yields
$\mathbb{E}[E_\tau] \le \mathbb{E}[M_\tau] \le 1$.
\end{proof}

We fit each factorisation with generalised Gaussian noise, whose shape parameter
spans Laplace through Gaussian to near-uniform, by profiling the shape over a
grid with iteratively reweighted least squares for the coefficients.

\begin{remark}[Abstention under Gaussianity]
When the noise is Gaussian the two factorisations are observationally
indistinguishable, the ratio~\eqref{eq:direction} does not grow, and no
certificate is issued. The procedure abstains precisely where the direction is
not identified rather than guessing. \citet{ding2025gaussdetect} give the
underlying equivalence: Gaussianity of the forward-model noise is equivalent to
independence between regressor and residual in the reverse regression.
\end{remark}

\begin{remark}[Two implementation constraints]
\label{rem:constraints}
Both were errors before they were understood. First, the observations scored by
the numerator and those entering the denominator's maximised likelihood must be
the \emph{same} index set; a denominator ranging over additional observations
subtracts their log-density with no matching numerator term, adding a constant
of order (extra observations) to $\log E_t$, which in our testing certified
\emph{both} directions of every pair, including under Gaussian noise. Second,
the conditioning block $Z$ is frozen when a pair's adjacency is certified and
the direction process accumulates only from that point; a block that kept
changing would change the hypothesis under test, whereas~\eqref{eq:ville} bounds
a fixed one.
\end{remark}

\begin{remark}[Misspecification]
\label{rem:misspec}
Proposition~\ref{prop:direction} needs the denominator's supremum to dominate
the likelihood at the true null parameter, which is not guaranteed when the
noise lies outside the fitted family. This is the certificate's main theoretical
risk. Empirically it is more robust than that suggests: against noise breaking
each family property in turn---Student-$t_3$ (polynomial rather than exponential
tails), shifted exponential (asymmetric) and a two-component mixture
(bimodal)---the reversed direction was certified in $0$ of $12$ runs in each
case at $\alpha = 0.05$, with full power retained.
\end{remark}

\section{CERT-CD}
\label{sec:algorithm}

The algorithm maintains a three-valued adjacency status for every pair. It
streams data in batches; after each batch it updates the running maximum of
$A_t(i,j)$ for every undecided pair and certifies adjacency when the maximum
crosses its threshold, at which point the two direction processes for that pair
are opened with a frozen conditioning block. The output graph contains only
certified edges; an edge carries an arrowhead when a direction certificate has
fired, and circles at both ends otherwise. \textbf{An absent edge means
undecided, not certified absent.}

Budgets are fixed before any data arrive: $\alpha = \alpha_A + \alpha_D$, with
$\alpha_A$ split evenly over the $\binom{d}{2}$ adjacency nulls and $\alpha_D$
over the $d(d-1)$ direction nulls.

\begin{theorem}[Whole-graph anytime validity]
\label{thm:main}
Let $\tau$ be any stopping time, including one chosen by inspecting the output.
Under the Markov condition, Assumption~\ref{ass:sep}, valid per-set e-processes,
and---for the direction certificates---the linear non-Gaussian model of
Section~\ref{sec:direction},
\begin{equation}
P\Bigl( \exists\, \tau : \text{ some certified edge is absent from } \Gstar
  \text{ or some certified arrowhead is wrong} \Bigr) \;\le\; \alpha .
\end{equation}
\end{theorem}

\begin{proof}
Each certificate is issued when an e-process for a null that holds whenever the
asserted feature is absent crosses $1/(\alpha_\bullet w_\bullet)$. By
Proposition~\ref{prop:adjacency}, Proposition~\ref{prop:direction} and Ville's
inequality~\eqref{eq:ville}, each is falsely issued with probability at most
$\alpha_\bullet w_\bullet$, uniformly in time. Both families are fixed before
any data are observed, so a union bound over the $\binom{d}{2}$ adjacency nulls
and $d(d-1)$ direction nulls gives $\alpha_A + \alpha_D = \alpha$. Which
certificates the algorithm evaluates never enters the bound, so the
graph-dependent order in which pairs are examined is free.
\end{proof}

\begin{remark}[On the multiplicity layer]
\label{rem:multiplicity}
The union bound is the least powerful admissible way to combine e-values;
e-closed testing with weighted e-Bonferroni local tests strictly dominates it
\citep{hartog2026fwer}. It cannot be used here. Our certificates reject on the
\emph{running maximum} of an e-process, and a running maximum is a
\emph{pseudo} e-value: Ville bounds the probability that it ever crosses, but
its expectation exceeds one, so it is not an e-value and closed testing does not
apply. \citet{hartog2026fwer} are explicit on this point, and recover a bound of
$\alpha + O(\alpha^2 \log \alpha^{-1})$ only for \emph{independent} e-processes.
Every per-hypothesis e-process here is a function of the same running Gram
matrix, so independence fails badly. The union bound is therefore the correct
instrument in this setting rather than a lazy one.
\end{remark}

\section{Related work}
\label{sec:related}

Most primitives used here are standard, and it is worth saying so precisely.

\paragraph{Not new.} That the minimum of e-values is an e-value for a union null
is standard \citep{vovk2021evalues}. Aggregating conditional-independence tests
over conditioning sets into an edge-level statistic is done by PC-p
\citep{strobl2017estimating}, which upper-bounds an edge's p-value by the
\emph{maximum} p-value over conditioning sets---the exact p-value analogue
of~\eqref{eq:adjacency}---and controls FDR with Benjamini--Yekutieli. Adjacency
as the object being tested is the premise of DAT \citep{amin2024scalable}, which
replaces the exponential set of tests with a relaxed optimisation problem solved
by neural networks. One-sided, faithfulness-free, finite-sample control of false
edges is exactly Theorem~2 of $\epsilon$-CUT \citep{shaska2025causal}, reached by
the same reasoning we use: a false-positive guarantee needs a correct null, not
faithfulness. Abstention and three-valued output are not new either: FCI leaves
unidentifiable edges undirected, \citet{uehara2026iterative} attaches per-edge
\texttt{resolved}/\texttt{impossible} codes with tiers that abstain when a
precondition test rejects, and \citet{prakash2024diagnostic} give a bivariate
direction test with an explicit inconclusive outcome. Sequential, anytime-valid
e-values \emph{inside} PC are already used by \citet{csillag2025prediction}, who
calibrate batched Fisher-$z$ p-values into per-test e-values. Residual
independence as a direction criterion is DirectLiNGAM
\citep{shimizu2011directlingam}, and \citet{ruiz2022sequentially} learn a
topological ordering from likelihood-ratio scores on residuals.

\paragraph{What is new.} Three things survive that accounting.

\emph{(i) Error-controlled orientation beyond the Markov equivalence class.}
Every method above with orientation error control is confined to the equivalence
class, so an edge in no v-structure remains unorientable; every method that
escapes the class decides directions by comparing scores, with no error control.
\citet{ruiz2022sequentially} is closest in statistic but proves only
population-level identifiability and requires all noise terms from one
scale-location family; its ``sequential'' refers to greedy sorting, not
sequential inference. \citet{uehara2026iterative} certifies \emph{provenance}---
which theorem licensed a direction---not error, and its Theorem~1 bounds expert
query counts under an ideal-oracle assumption.

\emph{(ii) A whole-graph guarantee proved rather than assumed.} Per-test
anytime-validity is not new, but \citet{csillag2025prediction} obtain graph-level
validity by assumption (their Assumption~2.5: the downstream algorithm is valid
whenever its inputs are), noting the multiple-comparison concern without
resolving it. Theorem~\ref{thm:main} proves the statement over a fixed family
and extends it to orientations.

\emph{(iii) Time-uniformity on the adjacency side.} $\epsilon$-CUT holds the
same one-sided, faithfulness-free ground \emph{at each fixed $n$}. Bounding the
probability that a false edge is \emph{ever} certified is a different statement,
and the one a monitored run needs; Section~\ref{sec:exp-calibration} measures the
gap. Secondarily, minimising over the complete family $\Scal_k$ performs all the
tests PC-p's maximum-p bound requires, so PC-p's zero-Type-II-error assumption
is not needed.

\section{Experiments}
\label{sec:experiments}

All experiments stream data in batches and inspect the estimate after every
batch, scoring the worst case over the whole trajectory rather than at a single
sample size. Code and figures accompany the paper.

\subsection{What monitoring costs a fixed-sample test}
\label{sec:exp-calibration}

Data are generated under a null conditional independence and the analyst
inspects the evidence as observations accumulate. Over 2000 replications
monitored to $n = 1000$, at nominal $\alpha = 0.05$ the Fisher-$z$ test rejects
in $34.6\%$ of runs when monitored, against $5.6\%$ evaluated once at the final
sample size; the e-process stays at $2.4\%$. This is the gap that separates a
fixed-$n$ guarantee from a time-uniform one.

The e-process pays for this in power. Against a fixed-sample Fisher-$z$ test at
the same $n$, evaluated once, it detects a coefficient of $0.08$ in $32\%$ of
runs against $72\%$, and $0.10$ in $58\%$ against $91\%$; the gap closes by
$0.15$ ($96\%$ against $100\%$). When the sample size can honestly be fixed in
advance, a fixed-sample test is the better instrument.

\subsection{Stratifying by the premise}

The competing guarantee for delete-on-non-rejection requires
$\delta$-strong faithfulness. This is a demanding condition: at $\delta = 0.15$
it holds in $65\%$ of random 6-variable DAGs and in only $19\%$ of 7-variable
graphs with one latent variable marginalised out, where the median smallest
partial correlation over true adjacencies is about $0.03$. Since removing an
edge whose association genuinely lies inside the equivalence region is correct
behaviour by construction, pooling both kinds of instance measures how often
random graphs are near-unfaithful rather than whether a procedure is calibrated.
We therefore stratify. Because the margin is a population quantity it can be
computed before simulating data, making it cheap to rejection-sample a
well-powered sample of each stratum.

\subsection{Certificate-only against delete-on-non-rejection}

Table~\ref{tab:main} reports the probability of ever making an unsound assertion
over a monitored run, at $\alpha = 0.1$, $d = 5$, Laplace noise.

\begin{table}[t]
\centering
\caption{Probability of an unsound assertion at any point in a monitored run
($\alpha = 0.1$). ``SPRINT-CD'' is an anytime-valid PC variant that deletes
edges against confidence-sequence certificates; both it and PC decide absence by
failing to reject.}
\label{tab:main}
\begin{tabular}{lcccc}
\toprule
& \multicolumn{2}{c}{CERT-CD} & SPRINT-CD & PC \\
\cmidrule(lr){2-3}
stratum & false edge & wrong arrow & lost edge & lost edge \\
\midrule
$\delta$-strong faithfulness holds ($n=40$) & 0.000 & 0.000 & 0.000 & 0.000 \\
$\delta$-strong faithfulness fails ($n=20$) & 0.050 & 0.000 & 0.600 & 0.800 \\
\bottomrule
\end{tabular}
\end{table}

Where the premise fails---the regime that breaks delete-on-non-rejection---
CERT-CD stays inside its budget while the two comparators fail in $60\%$ and
$80\%$ of runs. Near-unfaithful instances cost CERT-CD power: the affected pairs
stay undecided rather than being asserted wrongly.

On premise-satisfying instances at $n = 3000$, every true edge is certified and
$96\%$ are also oriented, against a \emph{Markov-equivalence-class orientation
ceiling of $33\%$}---the fraction a CPDAG can orient at all on these graphs.
That gap is the practical content of the direction certificate.

The same stratification applied to the skeleton alone ($d = 6$, 120
premise-satisfying instances) gives a probability of $0.008$ that a true edge is
ever lost, against $0.133$ for PC re-run at each sample size on the same
instances.

\subsection{Latent confounders}

An anytime-valid FCI variant using the same machinery recovers the oracle PAG
under one marginalised latent variable: on premise-satisfying instances
($n=80$) no oracle adjacency was ever lost at $\alpha = 0.1$, structural
Hamming distance to the oracle PAG fell from $12.7$ to $0.35$ by $n = 4000$, and
the bi-directed edge marking the latent was recovered in $7$ of $7$ cases where
one was present.

\section{Discussion}
\label{sec:discussion}

\paragraph{What the output does not say.} The absence of an edge means
undecided. CERT-CD cannot report that an edge is missing, and a pair may remain
undecided indefinitely if adjacency-faithfulness fails for it---which is exactly
the situation in which every delete-on-non-rejection method returns a confident
wrong answer. This is the price of the rule, and we regard reporting it as
preferable to concealing it.

\paragraph{Limitations.} The direction certificate assumes causal sufficiency
for the pair and a parametric noise family; Remark~\ref{rem:misspec} reports
what misspecification does empirically, but validity outside the family is not
established. The FCI variant implements orientation rules R1--R4 and omits
R5--R10, so its PAG is sound but not maximally informative. The direction
process refits per pair per batch and has been evaluated only up to $d = 7$;
scaling is unproven. Finally, the orientation results use Laplace noise at
$d = 5$, and the misspecification study of Remark~\ref{rem:misspec} is bivariate
rather than graph-level.

\paragraph{Terminology.} ``Anytime FCI'' \citep{spirtes2001anytime} denotes a
\emph{computational} property---the outer loop over conditioning-set sizes may be
interrupted and still return a sound graph. ``Anytime-valid'' here is an
\emph{inferential} property: time-uniform error control at data-dependent
stopping times. The two are unrelated.

\section*{Data and code availability}
Implementation, experiments and figures are available at
\url{https://github.com/vinhqdang/SPRINT-CD}.

\bibliographystyle{plainnat}
\begin{thebibliography}{99}

\bibitem[Amin and Wilson(2024)]{amin2024scalable}
A.~N. Amin and A.~G. Wilson.
\newblock Scalable and flexible causal discovery with an efficient test for adjacency.
\newblock In \emph{Proceedings of the 41st International Conference on Machine
  Learning}, PMLR 235, 2024.

\bibitem[Csillag et~al.(2025)]{csillag2025prediction}
D.~Csillag, C.~J. Struchiner, and G.~T. Goedert.
\newblock Prediction-powered e-values.
\newblock In \emph{Proceedings of the 42nd International Conference on Machine
  Learning}, 2025. arXiv:2502.04294.

\bibitem[Ding and Zhang(2025)]{ding2025gaussdetect}
Z.~Ding and X.-P. Zhang.
\newblock {GaussDetect-LiNGAM}: Causal direction identification without
  {G}aussianity test.
\newblock Preprint, December 2025. arXiv:2512.03428.

\bibitem[Gr\"unwald et~al.(2024)]{grunwald2024safe}
P.~Gr\"unwald, R.~de~Heide, and W.~Koolen.
\newblock Safe testing.
\newblock \emph{Journal of the Royal Statistical Society: Series B}, 86(5):1091--1128, 2024.

\bibitem[Hartog and Lei(2026)]{hartog2026fwer}
W.~Hartog and L.~Lei.
\newblock Family-wise error rate control with e-values.
\newblock Preprint, 2026. arXiv:2501.09015.

\bibitem[He and Sutherland(2026)]{he2026sequential}
Z.~He and D.~J. Sutherland.
\newblock Sequential kernel-based conditional independence testing via adaptive betting.
\newblock In \emph{Proceedings of the 43rd International Conference on Machine
  Learning}, 2026. arXiv:2606.18993.

\bibitem[Prakash et~al.(2024)]{prakash2024diagnostic}
S.~Prakash, F.~Xia, and E.~Erosheva.
\newblock A diagnostic tool for functional causal discovery.
\newblock Preprint, 2024. arXiv:2406.07787.

\bibitem[Ramdas et~al.(2023)]{ramdas2023gametheoretic}
A.~Ramdas, P.~Gr\"unwald, V.~Vovk, and G.~Shafer.
\newblock Game-theoretic statistics and safe anytime-valid inference.
\newblock \emph{Statistical Science}, 38(4):576--601, 2023.

\bibitem[Ramsey et~al.(2006)]{ramsey2006adjacency}
J.~Ramsey, P.~Spirtes, and J.~Zhang.
\newblock Adjacency-faithfulness and conservative causal inference.
\newblock In \emph{Proceedings of the 22nd Conference on Uncertainty in
  Artificial Intelligence}, pages 401--408, 2006.

\bibitem[Ruiz et~al.(2022)]{ruiz2022sequentially}
G.~Ruiz, O.~H. Madrid~Padilla, and Q.~Zhou.
\newblock Sequentially learning the topological ordering of causal directed
  acyclic graphs with likelihood ratio scores.
\newblock Preprint, 2022. arXiv:2202.01748.

\bibitem[Shaska and Mitra(2025)]{shaska2025causal}
J.~Shaska and U.~Mitra.
\newblock Causal link discovery with unequal edge error tolerance.
\newblock Preprint, 2025. arXiv:2507.21570.

\bibitem[Shimizu et~al.(2006)]{shimizu2006linear}
S.~Shimizu, P.~O. Hoyer, A.~Hyv\"arinen, and A.~Kerminen.
\newblock A linear non-{G}aussian acyclic model for causal discovery.
\newblock \emph{Journal of Machine Learning Research}, 7:2003--2030, 2006.

\bibitem[Shimizu et~al.(2011)]{shimizu2011directlingam}
S.~Shimizu, T.~Inazumi, Y.~Sogawa, A.~Hyv\"arinen, Y.~Kawahara, T.~Washio,
  P.~O. Hoyer, and K.~Bollen.
\newblock {DirectLiNGAM}: A direct method for learning a linear non-{G}aussian
  structural equation model.
\newblock \emph{Journal of Machine Learning Research}, 12:1225--1248, 2011.

\bibitem[Spirtes(2001)]{spirtes2001anytime}
P.~Spirtes.
\newblock An anytime algorithm for causal inference.
\newblock In \emph{Proceedings of the 8th International Workshop on Artificial
  Intelligence and Statistics}, pages 213--221, 2001.

\bibitem[Spirtes et~al.(2000)]{spirtes2000causation}
P.~Spirtes, C.~Glymour, and R.~Scheines.
\newblock \emph{Causation, Prediction, and Search}.
\newblock MIT Press, 2nd edition, 2000.

\bibitem[Strobl et~al.(2017)]{strobl2017estimating}
E.~V. Strobl, P.~L. Spirtes, and S.~Visweswaran.
\newblock Estimating and controlling the false discovery rate of the {PC}
  algorithm using edge-specific p-values.
\newblock Preprint, 2017. arXiv:1607.03975.

\bibitem[Uehara(2026)]{uehara2026iterative}
E.~Uehara.
\newblock Iterative causal discovery: Per-edge impossibility certificates,
  tier-aware oracle queries, and the $1+K$ lower bound.
\newblock Preprint, 2026. arXiv:2605.27477.

\bibitem[Vovk and Wang(2021)]{vovk2021evalues}
V.~Vovk and R.~Wang.
\newblock E-values: Calibration, combination and applications.
\newblock \emph{Annals of Statistics}, 49(3):1736--1754, 2021.

\bibitem[Wasserman et~al.(2020)]{wasserman2020universal}
L.~Wasserman, A.~Ramdas, and S.~Balakrishnan.
\newblock Universal inference.
\newblock \emph{Proceedings of the National Academy of Sciences},
  117(29):16880--16890, 2020.

\end{thebibliography}

\end{document}
